\begin{textsample}{POS Dim 3 – mg – Score 12.00 – mg/mg\_inf\_principios\_5.txt}  \label{ex:f3_pos_248}
2.4.3 Concepção de Educação \textbf{Infantil} No cenário educacional atual, o campo da Educação \textbf{Infantil} tem sido tema de constantes debates e discussões sobre as concepções que subsidiam as práticas pedagógicas mediadoras de aprendizagem e do desenvolvimento dos \textbf{bebês}, das \textbf{crianças} bem \textbf{pequenas} e \textbf{crianças} \textbf{pequenas} nas \textbf{creches} e \textbf{pré-escolas}.
De acordo com as DCNEI ( 2010, p.12 ) define -se a Educação \textbf{Infantil} como : > primeira etapa da educação básica, oferecida em \textbf{creches} e \textbf{pré-escolas}, às quais se caracterizam como espaços institucionais não \textbf{domésticos} que constituem estabelecimentos educacionais públicos ou privados que educam e cuidam de \textbf{crianças} de 0 a 5 anos de idade no período diurno, em jornada integral ou parcial, regulados e supervisionados por órgão competente do sistema de ensino e submetidos a controle social.
Nesta perspectiva, as instituições de Educação \textbf{Infantil} possuem um compromisso político e social ao garantir as especificidades da infância na sociedade atual ; ocupando um lugar importante e pensado para a \textbf{criança}, garantindo -lhe o seu papel central no planejamento em relação ao currículo e às práticas pedagógicas, pois quando não é assegurado o caráter ativo da \textbf{criança} no processo, a aprendizagem será sempre incompleta. 2.4.4 Ser Professor ( a ) na Educação \textbf{Infantil} A intencionalidade da ação docente constitui -se elemento essencial da prática pedagógica para a concretização de um trabalho comprometido com a garantia dos direitos de aprendizagem e desenvolvimento.
Nesse sentido, pressupõe -se que o ( a ) professor ( a ) desenvolva uma prática pedagógica reflexiva, pautada na \textbf{escuta}, no olhar sensível, na mediação e na promoção de um ambiente acolhedor e propício à aprendizagem e ao desenvolvimento. > O compromisso dos/as professores/as e das instituições de Educação \textbf{Infantil} é observar e interagir com as \textbf{crianças} e seus modos de expressar e elaborar saberes.
Com base nesse processo dinâmico de acolhimento dos saberes \textbf{infantis}, está à ação dos/as docentes em selecionar, organizar, refletir, mediar e avaliar o conjunto das práticas cotidianas que se realizam na escola, com a participação das \textbf{crianças}.
A partir disso, o/a professor/a promove interações das \textbf{crianças} com conhecimentos que fazem parte do patrimônio cultural, artístico, ambiental, científico e tecnológico, por meio do planejamento de possibilidades e oportunidades que se constituem a partir da observação, dos questionamentos e do diálogo constante com as \textbf{crianças} ( BRASIL, 2016, p. 59-60 ).
A atuação docente na Educação \textbf{Infantil} implica ter clareza das particularidades do trabalho com \textbf{bebês}, \textbf{crianças} bem \textbf{pequenas} e \textbf{crianças} \textbf{pequenas}, uma vez que exige um conhecimento específico sobre a aprendizagem, o desenvolvimento e as necessidades básicas ( cognitivas, sociais, emocionais, físicas ) da \textbf{criança}, para que possa promover práticas pedagógicas adequadas. > Há uma especificidade clara no trabalho do professor de Educação \textbf{Infantil} que é a de ter a sensibilidade para as linguagens da \textbf{criança}, para o estímulo à autonomia, para mediar a construção de conhecimentos científicos, artísticos e tecnológicos, e, também, para se colocar no lugar do outro, aspectos imprescindíveis no estabelecimento de vínculos com \textbf{bebês} e \textbf{crianças} \textbf{pequenas} ( BARBOSA, 2009, p.37 ).
As ações de cuidar e educar são complementares e indissociáveis, e constituem os pilares que sustentam as ações pedagógicas realizadas em \textbf{creches} e \textbf{pré-escolas}.
Tais elementos fundamentais da prática docente sinalizam que esta é uma tarefa complexa, que deve ser concretizada nas situações cotidianas de interações e \textbf{brincadeiras} que compõem a sua proposta curricular.
Entende -se que para o professor integrar o cuidar e o educar às suas práticas, deve -se partir da vinculação entre o saber e o saber fazer, uma vez que as necessidades apresentadas no contexto da Educação \textbf{Infantil} não permitem que uma dimensão anteceda a outra.
Nesta perspectiva, a formação continuada de professores é fundamental e deve ser conduzida a partir de um processo contínuo de reflexão em torno da prática profissional, que favoreça a construção de novos saberes. 2.4.5 A família e a escola Sendo a família a primeira educadora da \textbf{criança}, entende -se que é nela que a \textbf{criança} inicia as suas relações afetivas, encontra o outro, e por meio dele, aprende os modos de existir.
Dessa forma, o seu mundo adquire significado e ela começa a constituir -se como sujeito a partir dessa interação familiar.
A escola, por sua vez, possui papel complementar à ação da família e da comunidade no alcance do desenvolvimento integral da \textbf{criança}.
Isso implica, entre outras exigências, a necessidade de \textbf{acolher} as famílias, conhecer suas expectativas sobre a educação a ser ofertada a seus filhos, seu modo de ser e estar no mundo.
Além disso, em uma ação complementar à das famílias, da comunidade e do poder público, é imprescindível, como dispõe o Estatuto da Criança e do Adolescente — ECA que “ as instituições educativas assegurem o direito das \textbf{crianças} ao \textbf{cuidado}, à proteção, à saúde, à liberdade, à confiança, ao respeito, à dignidade, à cultura, às artes, à \textbf{brincadeira}, à convivência e à interação com outros/as meninos/as ” ( BRASIL, 2016, p. 55 ).
Sendo assim, faz -se necessário investir nessa integração e \textbf{apoiar} -se nos diversos benefícios da proximidade entre a família e a escola e, a partir disso, construir uma relação fundamentada em ações colaborativas e integradas, visando oferecer aos \textbf{bebês}, \textbf{crianças} bem \textbf{pequenas} e \textbf{crianças} \textbf{pequenas}, as melhores condições de desenvolvimento e aprendizagem.

% matched lemmas: acolher, apoiar, bebê, brincadeira, creche, criança, cuidado, doméstico, escuta, infantil, pequeno, pré-escola
\end{textsample}
